\chapter{Connections: Groups, Tensors, and Functional Analysis}

In Chapter~36, we used quadratic forms to find the lowest points of energy bowls,
completing our tour of finite-dimensional optimization. But linear algebra does
not end where the dimension stops being finite, or where matrices act on more than
two vectors.

You have come a long way. You started with arrows in the plane and the simple act
of adding them tip to tail. You learned to span, to project, to factor, to find
the directions a matrix refuses to turn. By now linear algebra probably feels
less like a subject and more like a \emph{language} -- a way of talking about
straightness, about combining things proportionally, about the geometry hiding
inside arrays of numbers.

This last chapter is a set of open doors. We are not going to master any of the
three subjects in the title; each is a whole field with its own thick books and
its own community. Our goal is to examine how each area builds on the foundational
concepts developed throughout this book. When you later encounter Lie groups,
quantum mechanics, the finite element method, or modern neural networks, you will
recognize these advanced frameworks as natural extensions of linear algebra.

Three doors. Behind the first, \textbf{groups}: the matrices that preserve
length and angle -- rotations and reflections -- form a structure of their own,
and the study of how such structures \emph{act} on vector spaces is called
representation theory, the bridge from linear algebra to symmetry. Behind the
second, \textbf{tensors}: what happens when ``linear in one vector'' becomes
``linear in several vectors at once,'' the natural home of multilinear maps, of
the determinant, and of every modern machine-learning framework. Behind the third,
\textbf{functional analysis}: what survives when the vectors become \emph{functions}
and the dimension becomes infinite, where Fourier series turn out to be nothing
more exotic than an orthonormal basis you cannot see all of at once.

\begin{focusbox}
The single idea threading this whole chapter: the patterns of finite-dimensional
linear algebra -- basis, orthogonality, projection, eigenvalues, the action of a
linear map -- do not stop at the edge of $\R^n$. They keep going. They organize
symmetry (groups), they stack up (tensors), and they survive the leap to infinity
(Hilbert spaces). These directions show how the core ideas of linear algebra extend
into broader mathematical settings.
\end{focusbox}

\section{Matrix groups: the geometry that linear algebra preserves}
\label{ch37:groups}

Let us begin where your geometric instincts are strongest: rigid motions of the
plane. Pick up the plane $\R^2$ and move it so that no distances change and no
angles change -- so that the unit square stays a unit square, just possibly turned
or flipped. What are you allowed to do? You may \emph{rotate} about the origin, and
you may \emph{reflect} across a line through the origin. That is all (if we insist
the origin stays put). Every such motion is a linear map, and we have met its
matrices before.

A rotation by angle $\theta$ has matrix
\[
  \mat{R}(\theta)=\begin{bmatrix}\cos\theta & -\sin\theta\\ \sin\theta & \cos\theta\end{bmatrix},
\]
and a reflection across the line at angle $\phi/2$ has matrix
\[
  \mat{F}(\phi)=\begin{bmatrix}\cos\phi & \sin\phi\\ \sin\phi & -\cos\phi\end{bmatrix}.
\]
What these matrices have in common is the defining property of \emph{not distorting}:
their columns are orthonormal. Recall from the orthogonality chapter that a matrix
with orthonormal columns satisfies $\mat{Q}\T\mat{Q}=\mat{I}$, equivalently
$\mat{Q}\T=\mat{Q}\inv$. Such matrices are called \emph{orthogonal}.

\begin{definition}[Orthogonal group]
\label{ch37:def-on}
The set of all $n\times n$ real matrices with $\mat{Q}\T\mat{Q}=\mat{I}$ is called
the \emph{orthogonal group} $O(n)$. The subset with $\det\mat{Q}=+1$ is the
\emph{special orthogonal group} $SO(n)$ -- the pure rotations.
\end{definition}

\begin{recoverybox}
\textbf{Pause and think.} Why does having orthonormal columns automatically force $\mat{Q}\T\mat{Q}=\mat{I}$? Write $\mat{Q} = [\vec{q}_1\,\cdots\,\vec{q}_n]$. Row $i$ of $\mat{Q}\T$ meets column $j$ of $\mat{Q}$ in the dot product $\vec{q}_i\T\vec{q}_j$. If $i=j$, the column is a unit vector, giving length-squared $1$. If $i\ne j$, the distinct columns are perpendicular, giving $0$. The product matrix has $1$s on the diagonal and $0$s everywhere else -- exactly the identity matrix $\mat{I}$.
\end{recoverybox}

Why the word \emph{group}? Because this collection of matrices is closed under
exactly the operations that matter for ``doing one motion after another.''

\begin{intuitionbox}
Forget the formal axioms for a moment and think physically. If you rotate the plane
and then rotate again, you have rotated -- a single motion of the same kind. If you
do a rigid motion, you can always \emph{undo} it with another rigid motion. And
``do nothing'' is itself a (boring) rigid motion. So the rigid motions form a
self-contained world: combine two and you stay inside, every one has an inverse
inside, and there is a do-nothing element. That self-containedness \emph{is} what
``group'' means.
\end{intuitionbox}

\begin{formalbox}
A \emph{group} is a set $G$ with a way to combine two elements (here, matrix
multiplication) such that: (1) the combination of two elements of $G$ is again in
$G$ (\emph{closure}); (2) the operation is associative; (3) there is an
\emph{identity} element $e$ with $e g = g e = g$ for all $g$ (here $\mat{I}$); and
(4) every $g$ has an \emph{inverse} $g\inv \in G$ with $g g\inv = e$ (here
$\mat{Q}\inv = \mat{Q}\T$). A group need not be commutative, and $O(n)$ for
$n\ge 2$ is not.
\end{formalbox}

\begin{historybox}
In the 1870s, Norwegian mathematician Sophus Lie set out to do for differential equations what \'Evariste Galois had done for polynomial equations: classify their solutions by their underlying continuous symmetries. Rather than relying on isolated integration tricks, Lie studied smooth families of geometric transformations---now called \textbf{Lie groups}---and discovered that their curved geometric structure is governed by their tangent directions at the do-nothing identity element. This tangent space is the \textbf{Lie algebra} of infinitesimal generators, which records instantaneous velocities of transformation (such as angular velocity during a rotation). Exponentiating a skew-symmetric matrix generator $e^{t\mat{X}}$ turns these linear infinitesimal steps into finite, smooth rotations in $SO(n)$, creating an exact bridge between linear algebra and differential equations.

In the late 1920s, Hermann Weyl and Eugene Wigner carried group theory into quantum physics, proving that physical symmetries dictate the existence of quantum spin, angular momentum, and the classification of elementary particles. Yet this abstraction met fierce controversy: prominent physicists initially revolted against what they called the ``\emph{Gruppenpest}'' (the ``group pestilence''), complaining that abstract group representations were needlessly mathematical compared to direct wave calculations. The pestilence won: today the Standard Model of particle physics is formulated entirely in the language of Lie gauge groups, $U(1)\times SU(2)\times SU(3)$.

Meanwhile, David Hilbert in the early 1900s and Stefan Banach in the 1920s pushed linear algebra beyond finite dimensions into complete spaces of functions: \textbf{Banach spaces} (vector spaces with lengths, or norms, where every Cauchy sequence converges) and \textbf{Hilbert spaces} (which add an inner product to measure angles and perpendicularity). Early mathematicians resisted treating functions as vectors, wondering how an infinite continuum of values could ever behave like a Euclidean column of numbers. Hilbert proved that in the square-integrable function space $L^2$, the inner product $\int f(x)g(x)\,dx$ gives functions lengths, angles, and orthogonal projections identical to Euclidean space.

When Albert Einstein, Boris Podolsky, and Nathan Rosen challenged quantum mechanics in 1935 with the EPR paradox, dismissing quantum entanglement as ``spooky action at a distance,'' linear algebra provided the exact resolution through \textbf{tensor products}. In a composite state space $\mathcal{H}_A \otimes \mathcal{H}_B$, an entangled state vector $\vec{w}$ cannot be factored as a simple product $\vec{u} \otimes \vec{v}$. Non-local correlations are an intrinsic geometric property of tensor product spaces. Today, major questions remain active: proving whether quantum Yang--Mills gauge theory has a positive mass gap is a Clay Millennium Problem, while in multilinear algebra, simply computing the rank of a 3-way tensor is NP-hard.
\end{historybox}

Let us check that $O(n)$ really is a group, because the proof is three lines of
linear algebra you can already do, and it shows how the abstract axioms cash out
into facts about matrices.

\begin{proposition}[$O(n)$ is a group]
\label{ch37:prop-ongroup}
If $\mat{P}\T\mat{P}=\mat{I}$ and $\mat{Q}\T\mat{Q}=\mat{I}$, then $\mat{PQ}$ is
orthogonal, and $\mat{P}\inv$ is orthogonal.
\end{proposition}

\begin{proof}
For the product, $(\mat{PQ})\T(\mat{PQ}) = \mat{Q}\T\mat{P}\T\mat{P}\mat{Q}
= \mat{Q}\T\mat{I}\mat{Q} = \mat{Q}\T\mat{Q} = \mat{I}$. For the inverse, since
$\mat{P}\inv = \mat{P}\T$ we have $(\mat{P}\inv)\T \mat{P}\inv
= (\mat{P}\T)\T\mat{P}\T = \mat{P}\mat{P}\T = \mat{I}$, where the last equality
holds because a square matrix with $\mat{P}\T\mat{P}=\mat{I}$ also satisfies
$\mat{P}\mat{P}\T=\mat{I}$. The identity $\mat{I}$ is orthogonal, and matrix
multiplication is always associative. All four axioms hold.
\end{proof}

The determinant condition cuts the orthogonal group into two clean halves. An
orthogonal matrix has $\det = \pm 1$ (because $1 = \det\mat{I} = \det(\mat{Q}\T\mat{Q})
= (\det\mat{Q})^2$). The $+1$ half is the rotations, $SO(n)$. The $-1$ half is the
reflections (and rotation--reflection combinations). The two halves cannot be
deformed into one another without passing through a matrix of determinant $0$,
which is not invertible -- so a reflection is geometrically a different \emph{kind}
of motion from a rotation. You cannot continuously turn a left hand into a right
hand.

\begin{figure}[H]
\centering
\begin{tikzpicture}[scale=1.0,>=Latex]
  % rotation panel
  \begin{scope}[xshift=0cm]
    \draw[help lines,gray!25] (-1.4,-1.4) grid (1.4,1.4);
    \draw[->] (-1.5,0)--(1.6,0); \draw[->] (0,-1.5)--(0,1.6);
    \fill[vecblue!15] (0,0)--(1,0)--(1,1)--(0,1)--cycle;
    \draw[vecblue,thick] (0,0)--(1,0)--(1,1)--(0,1)--cycle;
    \node[vecblue] at (0.75,0.22) {\small F};
    \begin{scope}[rotate=35]
      \fill[vecred!15] (0,0)--(1,0)--(1,1)--(0,1)--cycle;
      \draw[vecred,thick] (0,0)--(1,0)--(1,1)--(0,1)--cycle;
      \node[vecred] at (0.5,0.55) {\small F};
    \end{scope}
    \node[below] at (0,-1.7) {\small $SO(2)$: rotate, $\det=+1$};
  \end{scope}
  % reflection panel
  \begin{scope}[xshift=5.2cm]
    \draw[help lines,gray!25] (-1.4,-1.4) grid (1.4,1.4);
    \draw[->] (-1.5,0)--(1.6,0); \draw[->] (0,-1.5)--(0,1.6);
    \fill[vecblue!15] (0,0)--(1,0)--(1,1)--(0,1)--cycle;
    \draw[vecblue,thick] (0,0)--(1,0)--(1,1)--(0,1)--cycle;
    \draw[gray,dashed] (-1.4,-1.4)--(1.4,1.4);
    \node[vecblue] at (0.32,0.72) {\small F};
    \node[vecred] at (0.72,0.32) {\rotatebox{-90}{\reflectbox{\small F}}};
    \node[below] at (0,-1.7) {\small reflection, $\det=-1$};
  \end{scope}
\end{tikzpicture}
\caption{Left: a rotation in $SO(2)$ turns the unit square (and the letter F)
without flipping it. Right: a reflection across the dashed line $y=x$ mirrors the F
without distortion. Both preserve all lengths and angles, so both are orthogonal;
they live in the two disconnected halves $\det=+1$ and $\det=-1$ of $O(2)$.}
\label{ch37:fig-on}
\end{figure}

\subsection{$SO(2)$ is a circle, and that circle is the key to everything}
\label{ch37:so2}

Here is the first genuinely new idea. The rotations of the plane are parametrized
by a single angle $\theta$, and $\theta$ lives on a \emph{circle} (because
$\theta$ and $\theta + 2\pi$ give the same rotation). So $SO(2)$ is not just a set
of matrices -- it has a \emph{shape}, and that shape is a circle. A group that is
also a smooth geometric space (a \emph{manifold} -- a space that locally looks like
flat Euclidean space, allowing calculus to be done on it) is called a \emph{Lie group},
and from here the subject opens out into a vast territory of mathematics and physics.
Figure~\ref{ch37:fig-so2circle} draws the identification: each rotation matrix
$\mat{R}(\theta)$ is one point on the unit circle, and composing rotations slides
you around it by adding angles.

\begin{figure}[H]
\centering
\begin{tikzpicture}[scale=1.7,>=Latex]
  \draw[help lines,gray!25] (-1.25,-1.25) grid (1.25,1.25);
  \draw[->] (-1.4,0)--(1.5,0) node[right]{\small $\cos\theta$};
  \draw[->] (0,-1.4)--(0,1.5) node[above]{\small $\sin\theta$};
  \draw[structurepurple,thick] (0,0) circle (1);
  % a point at angle alpha
  \def\al{30}
  \def\be{65}
  \fill[vecblue] (\al:1) circle (1.2pt);
  \draw[vecblue,->] (0,0)--(\al:1);
  \node[vecblue,right=4pt] at (\al:1) {\small $\mat{R}(\alpha)$};
  \fill[vecred] (\be:1) circle (1.2pt);
  \node[vecred,above=3pt] at (\be:1) {\small $\mat{R}(\alpha{+}\beta)$};
  % angle arc from alpha to alpha+beta inside circle
  \draw[targetgreen,->,thick] (\al:0.65) arc (\al:\be:0.65);
  \node[targetgreen] at (47.5:0.80) {\small $+\beta$};
  % theta arc at origin
  \draw[vecblue] (0.28,0) arc (0:\al:0.28);
  \node[vecblue] at (15:0.42) {\small $\alpha$};
\end{tikzpicture}
\caption{$SO(2)$ \emph{is} a circle. Each rotation $\mat{R}(\theta)$ corresponds to
the point $(\cos\theta,\sin\theta)$, and the group law
$\mat{R}(\alpha)\mat{R}(\beta)=\mat{R}(\alpha+\beta)$ is just sliding along the
circle by the added angle $\beta$. ``Do nothing'' sits at angle $0$; the inverse of
a point is its mirror across the horizontal axis.}
\label{ch37:fig-so2circle}
\end{figure}

Watch how the group law becomes the addition of angles. Multiply two rotations:
\[
  \mat{R}(\alpha)\mat{R}(\beta)
  = \begin{bmatrix}\cos\alpha & -\sin\alpha\\ \sin\alpha & \cos\alpha\end{bmatrix}
    \begin{bmatrix}\cos\beta & -\sin\beta\\ \sin\beta & \cos\beta\end{bmatrix}
  = \begin{bmatrix}\cos(\alpha+\beta) & -\sin(\alpha+\beta)\\ \sin(\alpha+\beta) & \cos(\alpha+\beta)\end{bmatrix}
  = \mat{R}(\alpha+\beta).
\]
Look closely: working out that product \emph{is} a derivation of the angle-addition
formulas for sine and cosine. The top-left entry is $\cos\alpha\cos\beta -
\sin\alpha\sin\beta$, which equals $\cos(\alpha+\beta)$. Trigonometry's hardest
identities fall out of the statement ``rotating by $\alpha$ then by $\beta$ is the
same as rotating by $\alpha+\beta$.'' The group law and the angle-addition identity
express the same algebraic relationship.

\begin{connectionbox}
This is why unit complex numbers and planar rotations have the exact same algebra. Identify
the matrix $\mat{R}(\theta)$ with the complex number $e^{i\theta} = \cos\theta +
i\sin\theta$, sitting on the unit circle in $\C$. Then $\mat{R}(\alpha)\mat{R}(\beta)
= \mat{R}(\alpha+\beta)$ becomes $e^{i\alpha}e^{i\beta} = e^{i(\alpha+\beta)}$, the
ordinary rule for multiplying exponentials. The circle group $SO(2)$, the unit
complex numbers $\{e^{i\theta}\}$, and the angles modulo $2\pi$ are three faces of
one object. When you later meet ``$U(1)$ gauge symmetry'' in physics, it is this
circle, and nothing more frightening.
\end{connectionbox}

In three dimensions the story gets richer and more physical. $SO(3)$, the rotations
of space, is a three-dimensional Lie group (you need three numbers to specify a
rotation: an axis, which costs two numbers, and an angle of turn, which costs one).
Euler's theorem -- itself pure linear algebra -- says every element of $SO(3)$ is a
rotation about \emph{some} axis. We can prove this directly from the eigenvalues of $\mat{R}$.

\begin{proposition}[Every rotation of space has an axis]
\label{ch37:prop-euler}
Every $\mat{R}\in SO(3)$ has $+1$ as an eigenvalue. The eigenvector is the axis of
rotation -- the line $\mat{R}$ leaves fixed.
\end{proposition}

\begin{proof}
The eigenvalues of $\mat{R}$ are the roots of its characteristic polynomial
$\det(\mat{R}-\lambda\mat{I})$, a real cubic, so at least one root is real (complex
roots of a real polynomial come in conjugate pairs). Being orthogonal forces every
eigenvalue to have modulus $1$ (an orthogonal map preserves length, so
$\norm{\mat{R}\vec{x}}=\norm{\vec{x}}$ forces $\abs{\lambda}=1$ on any eigenvector),
so the real eigenvalues are $\pm 1$. Their product is $\det\mat{R}=+1$. The three
eigenvalues are therefore
either $\{1,1,1\}$, or $\{1,-1,-1\}$, or $\{1, e^{i\theta}, e^{-i\theta}\}$ -- in
every case a $+1$ appears. The vector $\vec{x}$ with $\mat{R}\vec{x}=\vec{x}$ is a
fixed direction, i.e. the axis.
\end{proof}

This is why a spinning top has an axis, why the Earth has a North Pole, and why
any sequence of 3D rotations you compose is still just \emph{one} rotation about
\emph{one} axis.
Figure~\ref{ch37:fig-euleraxis} shows the picture: the $+1$-eigenvector is the
fixed line, and the plane perpendicular to it spins by some angle $\theta$.

\begin{figure}[H]
\centering
\begin{tikzpicture}[scale=1.0,>=Latex,line join=round]
  % axis (the +1 eigenvector)
  \draw[structurepurple,very thick,->] (0,-1.5)--(0,2.1) node[above]{\small axis: $\mat{R}\vec{x}=\vec{x}$};
  % perpendicular plane drawn as an ellipse (foreshortened disk)
  \draw[targetgreen!70,thick] (0,0) ellipse (2 and 0.7);
  \node[targetgreen!60!black] at (2.5,-0.45) {\small plane $\perp$ axis};
  % a vector in the plane and its rotated image
  \draw[vecblue,thick,->] (0,0)--(2,0);
  \node[vecblue,right] at (2,0.02) {\small $\vec{w}$};
  \draw[vecred,thick,->] (0,0)--(1.0,0.606);
  \node[vecred,above left=1pt] at (1.0,0.606) {\small $\mat{R}\vec{w}$};
  % rotation arc inside the disk
  \draw[->] (1.3,0.15) arc (15:50:1.3 and 0.45);
  \node at (1.45,0.45) {\small $\theta$};
  \fill[structurepurple] (0,0) circle (1.4pt);
\end{tikzpicture}
\caption{Euler's theorem geometrically. Every $\mat{R}\in SO(3)$ fixes a line (its
$+1$-eigenvector, the purple axis) and rotates the perpendicular plane by an angle
$\theta$; the complex eigenvalues $e^{\pm i\theta}$ record that in-plane spin. A
composition of rotations is again one such axis-plus-angle.}
\label{ch37:fig-euleraxis}
\end{figure}

\begin{applicationbox}
\textbf{Rotations in graphics, robotics, and aerospace.} Every time a game engine
turns your character's head, or a drone holds level, or a satellite points its
camera, a member of $SO(3)$ is being computed. Because composing rotations as
$3\times 3$ matrices accumulates numerical rounding error (the matrix slowly stops
being strictly orthogonal), practitioners often store rotations as \emph{unit quaternions} --
four-component numbers $q = w + x i + y j + z k$ with $\norm{q}=1$. Unit quaternions
avoid ``gimbal lock'' (a physical trap where two rotation axes align and lose a degree
of freedom) and form a \emph{double cover} of $SO(3)$, meaning that the two antipodal
quaternions $q$ and $-q$ correspond to the exact same spatial rotation. They renormalize
with a single vector scaling and interpolate smoothly along the shortest arc on the
four-dimensional sphere (an algorithm known as \emph{slerp}, for spherical linear
interpolation). That quaternions are the ideal representation is a geometric theorem
about the topology of $SO(3)$, discovered by studying its symmetry group.
\end{applicationbox}

\subsection{The complex cousin: unitary matrices}
\label{ch37:unitary}

Everything above has a complex twin. Replace ``preserves length in $\R^n$'' with
``preserves length in $\C^n$,'' where length uses the Hermitian inner product
$\inner{\vec{u}}{\vec{v}} = \sum_k \overline{u_k} v_k$. The condition
$\mat{Q}\T\mat{Q}=\mat{I}$ becomes $\mat{U}^{*}\mat{U}=\mat{I}$, where $\mat{U}^{*}$
is the conjugate transpose. Such matrices are \emph{unitary}, and they form the
\emph{unitary group} $U(n)$.

\begin{focusbox}
Orthogonal is to real space as unitary is to complex space. $O(n)$ preserves
$\norm{\vec{x}}^2 = \vec{x}\T\vec{x}$; $U(n)$ preserves $\norm{\vec{x}}^2 =
\vec{x}^{*}\vec{x}$. Both are the ``rigid motions'' of their space -- the matrices
that never stretch, shrink, or shear.
\end{focusbox}

The unitary group is the mathematical home of quantum mechanics: the state of a
quantum system is a unit vector in $\C^n$, and the laws of physics evolve it by
unitary matrices, precisely because unitary matrices are the ones that keep the
total probability (the squared length) equal to $1$. When you read that ``time
evolution is unitary,'' it means ``the dynamics is a rigid rotation of the complex
state space.'' The Fourier matrix from the complex chapter is unitary (after
scaling by $1/\sqrt{n}$); so is every quantum logic gate.

\subsection{Representation theory: letting a group act}
\label{ch37:reps}

We have been treating groups as collections of matrices. But many groups arise
abstractly -- the symmetries of a molecule, the permutations of a deck of cards,
the rotations that leave a crystal looking the same -- with no matrices in sight.
\emph{Representation theory} is the art of giving such a group a set of matrices to
\emph{be}, so that the full toolbox of linear algebra can be brought to bear on it.

\begin{definition}[Representation]
\label{ch37:def-rep}
A \emph{representation} of a group $G$ on a vector space $V$ is a rule
$\rho$ that assigns to each group element $g$ an invertible matrix $\rho(g)$ acting
on $V$, in a way that respects the group law:
$\rho(g h) = \rho(g)\,\rho(h)$.
\end{definition}

The slogan is: \emph{a representation turns multiplication in the group into matrix
multiplication}. In plain terms, if you have an abstract symmetry operation, a
representation gives it a concrete matrix so it can actually transform vectors.
Once you have that, you can ask all your favorite linear-algebra questions about the
symmetry. Which directions does every symmetry leave alone (the \emph{invariant
subspaces})? Can we change basis so that every $\rho(g)$ becomes block-diagonal at
once, splitting the space into atomic, indivisible building blocks (the
\emph{irreducible representations} -- representations that contain no smaller invariant
subspaces)? This last question -- simultaneously block-diagonalizing a whole family of
matrices -- is the deep generalization of diagonalizing a single matrix.

\begin{example}[The smallest interesting representation]
\label{ch37:ex-rep}
Take the group with two elements: ``do nothing'' ($e$) and ``flip'' ($f$), where
$f f = e$. Represent it on $\R^2$ by
\[
  \rho(e) = \begin{bmatrix}1&0\\0&1\end{bmatrix},\qquad
  \rho(f) = \begin{bmatrix}0&1\\1&0\end{bmatrix}.
\]
Check the group law: $\rho(f)\rho(f) = \begin{bmatrix}0&1\\1&0\end{bmatrix}^2 =
\begin{bmatrix}1&0\\0&1\end{bmatrix} = \rho(e)$. Good. Now ask the
representation-theory question: change basis to make \emph{both} matrices diagonal
at once. The flip $\rho(f)$ has eigenvectors $\colvec{1\\1}$ (eigenvalue $+1$) and
$\colvec{1\\-1}$ (eigenvalue $-1$); the identity is already diagonal in any basis.
In the basis $\{\colvec{1\\1},\colvec{1\\-1}\}$ \emph{both} group elements are
diagonal:
\[
  \rho(e)\rightsquigarrow\begin{bmatrix}1&0\\0&1\end{bmatrix},\qquad
  \rho(f)\rightsquigarrow\begin{bmatrix}1&0\\0&-1\end{bmatrix}.
\]
The two-dimensional representation has split into two one-dimensional pieces: the
\emph{symmetric} line $\colvec{1\\1}$, on which the flip does nothing, and the
\emph{antisymmetric} line $\colvec{1\\-1}$, on which the flip negates. Those are the
two irreducible representations of the flip group. Every even/odd decomposition you
have ever seen -- even and odd functions, symmetric and antisymmetric tensors -- is
this same split. Figure~\ref{ch37:fig-rep} shows the two invariant lines and the
coordinate-swap geometry behind them.
\end{example}

\begin{recoverybox}
\textbf{Pause and think.} Look at how any vector decomposes along these two invariant lines. Take $\vec{x}=\colvec{x_1\\x_2}$. We can split it into symmetric and antisymmetric halves:
\[
  \colvec{x_1\\x_2} = \underbrace{\frac{x_1+x_2}{2}\colvec{1\\1}}_{\text{symmetric ($+1$ eigenvector)}} + \underbrace{\frac{x_1-x_2}{2}\colvec{1\\-1}}_{\text{antisymmetric ($-1$ eigenvector)}}.
\]
Applying the flip swaps coordinates, which leaves the symmetric piece untouched and negates the antisymmetric piece. The group symmetry has decoupled $\R^2$ into two independent, uncoupled 1D worlds!
\end{recoverybox}

\begin{figure}[H]
\centering
\begin{tikzpicture}[scale=1.5,>=Latex]
  \draw[help lines,gray!25] (-1.3,-1.3) grid (1.3,1.3);
  \draw[->] (-1.4,0)--(1.5,0) node[right]{\small $x_1$};
  \draw[->] (0,-1.4)--(0,1.5) node[above]{\small $x_2$};
  % invariant lines: symmetric (1,1) and antisymmetric (1,-1)
  \draw[structurepurple,thick] (-1.25,-1.25)--(1.25,1.25);
  \draw[targetgreen!70!black,thick] (-1.25,1.25)--(1.25,-1.25);
  \node[structurepurple,right] at (1.25,1.25) {\small $\colvec{1\\1}$: flip fixes ($+1$)};
  \node[targetgreen!60!black,below right] at (0.9,-0.95) {\small $\colvec{1\\-1}$: flip negates ($-1$)};
  % a vector and its image under rho(f) (swap coordinates)
  \def\ax{0.9} \def\ay{0.3}
  \draw[vecblue,thick,->] (0,0)--(\ax,\ay);
  \node[vecblue,right] at (\ax,\ay) {\small $\vec{x}$};
  \draw[vecred,thick,->] (0,0)--(\ay,\ax);
  \node[vecred,above left] at (\ay,\ax) {\small $\rho(f)\vec{x}$};
  % mirror dashes showing reflection across purple line
  \draw[gray,dashed] (\ax,\ay)--(\ay,\ax);
\end{tikzpicture}
\caption{The flip $\rho(f)=\protect\begin{bmatrix}0&1\\1&0\protect\end{bmatrix}$ swaps
coordinates, which is reflection across the symmetric line $\protect\colvec{1\\1}$.
That line and the antisymmetric line $\protect\colvec{1\\-1}$ are the invariant
subspaces: in this rotated basis \emph{both} group matrices are diagonal at once.
Simultaneously diagonalizing the whole family is the representation-theory analogue
of diagonalizing a single matrix.}
\label{ch37:fig-rep}
\end{figure}

\begin{aestheticbox}
Representation theory is one of the most beautiful unifications in mathematics
because it explains \emph{why} eigenvectors and ``special bases'' keep appearing in
physics. A physical system with a symmetry can always be organized into the
irreducible pieces of that symmetry, and those pieces are forced to be the natural
coordinates. The energy levels of the hydrogen atom, the multiplets of particle
physics, the selection rules of spectroscopy, the normal modes of a symmetric
molecule -- all are the irreducible representations of the system's symmetry group,
found by simultaneously diagonalizing a family of matrices. Diagonalization, which
you learned to do for one matrix, becomes the engine that classifies the structure
of the physical world.
\end{aestheticbox}

\begin{strangbox}
Strang's course gives you orthogonal matrices, $\mat{Q}\T\mat{Q}=\mat{I}$, and the
spectral theorem for symmetric matrices, but it (rightly, for a first course) stops
at single matrices. The leap here is conceptual: collect the orthogonal matrices into
the \emph{group} $O(n)$, and ask not ``diagonalize this one matrix'' but ``find the
common structure of this whole family of matrices.'' That reframing -- from one
operator to a symmetry group of operators -- is the doorway from a linear-algebra
course into modern geometry and physics.
\end{strangbox}

\section{Tensors: linearity in several slots at once}
\label{ch37:tensors}

So far a linear map takes \emph{one} vector as input and acts linearly on it:
$\mat{A}(\alpha\vec{x}+\beta\vec{y}) = \alpha\mat{A}\vec{x}+\beta\mat{A}\vec{y}$.
What if a map takes \emph{two} vectors at once and is linear in \emph{each} slot
separately? That is the idea of a tensor, and you have been using one since
the dot-product chapter without naming it.

\subsection{The dot product is a bilinear form}
\label{ch37:bilinear}

Consider the map that takes two vectors $\vec{x},\vec{y}\in\R^n$ and returns the
number $\vec{x}\T\vec{y}$. Hold $\vec{y}$ fixed and it is linear in $\vec{x}$; hold
$\vec{x}$ fixed and it is linear in $\vec{y}$. Linear in each slot -- that is
\emph{bilinear}. More generally a bilinear form is
\[
  B(\vec{x},\vec{y}) = \vec{x}\T\mat{M}\vec{y}
  = \sum_{i,j} M_{ij}\, x_i\, y_j,
\]
and the matrix $\mat{M}$ \emph{is} the tensor: a grid of numbers $M_{ij}$, one for
each pair (row-slot, column-slot). The quadratic forms from the optimization
chapter are the special case $\vec{y}=\vec{x}$.

\begin{intuitionbox}
A vector is a list of numbers indexed by \emph{one} index: $v_i$. A matrix is a
grid indexed by \emph{two}: $M_{ij}$. A tensor just keeps going -- a box of numbers
indexed by three, $T_{ijk}$, or more. The number of indices is the \emph{order} (or
``rank,'' though that word is overloaded). Vectors are order-1 tensors, matrices
are order-2 tensors, and the determinant of a $3\times 3$ matrix is computed by an
order-3 tensor: the Levi-Civita symbol $\varepsilon_{ijk}$, a bookkeeping device
that equals $+1$ for even permutations of $(1,2,3)$, $-1$ for odd permutations, and
$0$ whenever any index repeats. The natural generalization is to allow any number
of indices.
\end{intuitionbox}

The key property that makes tensors more than ``arrays'' is how they transform when
you change basis. A genuine tensor carries one transformation rule per index, so
that the multilinear map it represents stays the same object even as its
coordinates change. This is the precise sense in which ``a tensor is a thing that
transforms like a tensor'' -- a slogan that is circular only until you have seen the
change-of-basis chapter, after which it is exactly right: a tensor is a multilinear
map, and its array of components is basis-dependent the same way a matrix's entries
are.

\subsection{The tensor product: building bigger spaces by multiplication}
\label{ch37:tensorproduct}

Here is the constructive heart of the subject. Given two vectors $\vec{u}\in\R^m$
and $\vec{v}\in\R^n$, form their \emph{outer product}, the $m\times n$ matrix
\[
  \vec{u}\otimes\vec{v} = \vec{u}\,\vec{v}\T,
  \qquad (\vec{u}\otimes\vec{v})_{ij} = u_i v_j.
\]
You met this in the SVD chapter: a rank-one matrix is exactly an outer product. The
collection of all such products, and all their sums, is the \emph{tensor product
space} $\R^m\otimes\R^n$, which can be identified with the space of $m\times n$ matrices.
Its dimension is $m\cdot n$ -- dimensions \emph{multiply} under
$\otimes$, where they would \emph{add} under the direct sum $\oplus$.

\begin{labbox}
Compute one outer product by hand and watch the rank-one structure appear. With
$\vec{u}=\colvec{1\\2}$ and $\vec{v}=\colvec{3\\1\\-1}$,
\[
  \vec{u}\otimes\vec{v} = \colvec{1\\2}\begin{bmatrix}3&1&-1\end{bmatrix}
  = \begin{bmatrix}3&1&-1\\ 6&2&-2\end{bmatrix}.
\]
Notice every row is a multiple of $\begin{bmatrix}3&1&-1\end{bmatrix}$ and every
column is a multiple of $\colvec{1\\2}$. That is what rank one \emph{means}, and it
is the simplest possible nonzero tensor. The SVD decomposes any matrix into a
linear combination of rank-one tensor products:
$\mat{A}=\sum_k \sigma_k\,\vec{u}_k\otimes\vec{v}_k$.
\end{labbox}

\begin{recoverybox}
\textbf{Pause and think.} Why is any nonzero outer product $\vec{u}\otimes\vec{v} = \vec{u}\vec{v}\T$ automatically rank 1? Look at column $j$: it is simply $v_j \vec{u}$. Every single column is a scalar multiple of $\vec{u}$, so all columns lie on the one-dimensional line $\spn\{\vec{u}\}$. The column space has dimension 1!
\end{recoverybox}

\begin{figure}[H]
\centering
\begin{tikzpicture}[scale=0.82,>=Latex]
  % Panel 1: Rank-1 Outer Product
  \begin{scope}[shift={(0,0)}]
    \node[anchor=south,font=\footnotesize\scshape,text=spacegray!90!black] at (2.2,2.8) {1. Rank-1 Outer Product $\vec{u}\otimes\vec{v}$};
    \draw[fill=codebg,draw=spacegray!30,rounded corners=4pt] (-0.8,-1.8) rectangle (5.2,2.6);

    % u column
    \node[draw=sourceblue!70,fill=sourceblue!15,rounded corners=2pt,font=\scriptsize,minimum width=0.8cm] at (0,1.2) {$1$};
    \node[draw=sourceblue!70,fill=sourceblue!15,rounded corners=2pt,font=\scriptsize,minimum width=0.8cm] at (0,0.4) {$2$};
    \node[font=\footnotesize,text=sourceblue!90!black] at (0,-0.4) {$\vec{u}$};

    \node[font=\large] at (0.9,0.8) {$\otimes$};

    % v^T row
    \node[draw=targetgreen!70,fill=targetgreen!15,rounded corners=2pt,font=\scriptsize,minimum height=0.6cm] at (1.8,1.9) {$3$};
    \node[draw=targetgreen!70,fill=targetgreen!15,rounded corners=2pt,font=\scriptsize,minimum height=0.6cm] at (2.6,1.9) {$1$};
    \node[draw=targetgreen!70,fill=targetgreen!15,rounded corners=2pt,font=\scriptsize,minimum height=0.6cm] at (3.4,1.9) {$-1$};
    \node[font=\footnotesize,text=targetgreen!90!black] at (4.3,1.9) {$\vec{v}\T$};

    % result grid
    \foreach \i/\u in {0/2,1/1} {
      \foreach \j/\v/\val in {0/3/6,1/1/2,2/{-1}/{-2}} {
        \ifnum\i=1
          \pgfmathtruncatemacro{\pval}{\u*\v}
          \node[draw=spacegray!50,fill=white,rounded corners=1pt,font=\scriptsize,minimum width=0.75cm,minimum height=0.6cm] at (1.8+0.8*\j,0.4+0.8*\i) {$\pval$};
        \else
          \node[draw=spacegray!50,fill=white,rounded corners=1pt,font=\scriptsize,minimum width=0.75cm,minimum height=0.6cm] at (1.8+0.8*\j,0.4+0.8*\i) {$\val$};
        \fi
      }
    }
    \node[draw=sourceblue!50,fill=sourceblue!8,rounded corners=2pt,font=\tiny,align=center] at (2.2,-1.2) {
      Every row proportional to $\vec{v}\T$;\\
      every col proportional to $\vec{u}$ (rank 1)
    };
  \end{scope}

  % Panel 2: General Tensor Space V \otimes W & Entanglement
  \begin{scope}[shift={(6.6,0)}]
    \node[anchor=south,font=\footnotesize\scshape,text=spacegray!90!black] at (2.4,2.8) {2. General Tensor Space $V\otimes W$};
    \draw[fill=codebg,draw=spacegray!30,rounded corners=4pt] (-0.6,-1.8) rectangle (5.4,2.6);

    % Basis grid of e_i \otimes f_j
    \draw[step=0.8cm,gray!30,very thin] (0.4,0.1) grid (3.6,1.7);
    \foreach \r in {1,2} {
      \foreach \c in {1,2,3,4} {
        \fill[vecblue] (0.4+0.8*\c-0.4,0.1+0.8*\r-0.4) circle (1.8pt);
      }
    }
    \node[font=\tiny,text=vecblue] at (1.2,1.55) {$\vec{e}_1\otimes\vec{f}_2$};
    \node[font=\tiny,text=vecblue] at (2.8,0.75) {$\vec{e}_2\otimes\vec{f}_4$};

    \node[draw=quotientorange!60,fill=white,rounded corners=2pt,font=\tiny,align=center] at (2.4,-1.2) {
      $\dim(V\otimes W) = (\dim V)(\dim W)$\\
      General tensor can be \textbf{entangled}\\
      (cannot factor as $\vec{a}\otimes\vec{b}$)
    };
  \end{scope}

  % Panel 3: Order-3 Tensor Hyper-matrix
  \begin{scope}[shift={(13.6,0)}]
    \node[anchor=south,font=\footnotesize\scshape,text=spacegray!90!black] at (1.4,2.8) {3. 3-Way Tensor Cube $\mathcal{T}$};
    \draw[fill=codebg,draw=spacegray!30,rounded corners=4pt] (-0.8,-1.8) rectangle (3.6,2.6);

    % 3D isometric box
    \draw[fill=structurepurple!15,draw=structurepurple!70] (0,0) rectangle (1.8,1.4);
    \draw[fill=structurepurple!25,draw=structurepurple!70] (0,1.4) -- (0.8,2.0) -- (2.6,2.0) -- (1.8,1.4) -- cycle;
    \draw[fill=structurepurple!35,draw=structurepurple!70] (1.8,0) -- (2.6,0.6) -- (2.6,2.0) -- (1.8,1.4) -- cycle;

    % Slice lines
    \draw[dashed,gray!60] (0.9,0) -- (0.9,0.5);
    \draw[dashed,gray!60] (0.9,0.9) -- (0.9,1.4) -- (1.7,2.0);
    \draw[dashed,gray!60] (0,0.7) -- (0.55,0.7);
    \draw[dashed,gray!60] (1.25,0.7) -- (1.8,0.7) -- (2.6,1.3);

    \node[draw=structurepurple!70,fill=white,rounded corners=1pt,font=\scriptsize,inner sep=1.5pt,text=structurepurple!90!black] at (0.9,0.7) {$\mathcal{T}_{ijk}$};
    \node[font=\tiny,below] at (0.9,-0.1) {Mode 1};
    \node[font=\tiny,right=3pt] at (2.6,1.0) {Mode 2};
    \node[font=\tiny,above] at (1.3,2.0) {Mode 3};

    \node[draw=structurepurple!50,fill=white,rounded corners=2pt,font=\tiny,align=center] at (1.4,-1.2) {
      Multimodal array $\R^{I \times J \times K}$\\
      slices, fibers, factorizations
    };
  \end{scope}
\end{tikzpicture}
\caption{The geometry of the tensor product. Left: A rank-one tensor is the outer product $\vec{u}\otimes\vec{v} = \vec{u}\vec{v}\T$, where entries are pairwise products and every column is proportional to $\vec{u}$. Middle: The full tensor space $V\otimes W$ is the linear span of all such outer products, with dimension multiplying: $\dim(V\otimes W) = (\dim V)(\dim W)$. High-rank tensors are \emph{entangled} (cannot be factored as a single outer product). Right: Higher-order tensors form multi-way arrays $\mathcal{T} \in \R^{I \times J \times K}$ indexed along multiple modes, analyzed via tensor factorizations (such as CP and Tucker decompositions, which generalize the SVD to multi-index data).}
\label{ch37:fig-outer}
\end{figure}

\subsection{The determinant is an alternating multilinear form}
\label{ch37:det-tensor}

You have already met one of the most important tensors in mathematics and not
realized its pedigree: the determinant. Think of $\det$ as a map that takes the
$n$ \emph{columns} of a matrix, $\vec{a}_1,\dots,\vec{a}_n$, and returns a number.
From the determinant chapter you know two of its properties:
\begin{itemize}
  \item it is linear in each column separately (scale a column, the determinant
        scales; add columns, the determinant adds) -- so $\det$ is
        \emph{multilinear};
  \item swapping two columns flips the sign -- so $\det$ is \emph{alternating}.
\end{itemize}
A multilinear, alternating function of $n$ vectors in $\R^n$ is unique up to scale,
and fixing $\det\mat{I}=1$ pins it down completely. That single sentence is the
modern definition of the determinant, and it explains the signed-volume picture
geometrically: the only consistent way to assign a signed volume to a parallelepiped
that is additive in each edge and flips sign when you swap two edges is the
determinant. The whole theory of $\Lambda^n$, the ``exterior algebra'' (the algebra
of oriented lengths, areas, and volumes) that underlies differential forms in
multivariable calculus, grows from this one alternating tensor.

\begin{connectionbox}
\textbf{Tensors in calculus and relativity.} The gradient is a tensor (it tells you
how a scalar changes per direction); the Jacobian is a tensor (linear-in-each-slot
change of variables); the metric tensor $g_{ij}$ of a curved space is a bilinear
form $\inner{\vec{u}}{\vec{v}}_g = \sum g_{ij}u^i v^j$ that measures length when the
usual dot product no longer applies. Einstein's general relativity is written
entirely in tensors because the laws of physics must be the same in every coordinate
system, and a tensor is precisely an object whose meaning survives a change of
basis. Differential forms, the objects you integrate in multivariable calculus, are
the alternating tensors of this section.
\end{connectionbox}

\begin{applicationbox}
\textbf{Why deep learning runs on tensors.} The libraries powering modern machine
learning -- PyTorch, TensorFlow (the name is a giveaway), JAX -- have one central
data type: the tensor, a multi-index array. A batch of $B$ color images of size
$H\times W$ is a four-index tensor $X_{b,c,h,w}$ (batch, channel, height, width).
A neural network is a long composition of multilinear operations (matrix
multiplies, convolutions, which are tensor contractions) interleaved with simple
nonlinearities. \emph{Tensor contraction} -- summing over a shared index, exactly
like matrix multiplication's $\sum_k A_{ik}B_{kj}$ -- is the central computational
operation, and the entire efficiency of GPUs comes from doing billions of these
contractions in parallel. When you hear ``GPT has a trillion parameters,'' those
parameters live in tensors, and training adjusts them by gradients that are
themselves computed by the chain rule applied to tensor contractions. The order-2
special case you already know, $\mat{A}\mat{B}$, is the seed of all of it.
\end{applicationbox}

\begin{trapbox}
``Tensor'' in a machine-learning library means only \emph{multidimensional array} --
a box of numbers with a shape. The mathematician's tensor carries the extra promise
of a change-of-basis transformation law. A PyTorch tensor of shape $(3,3)$ is not
automatically a geometric tensor in the relativity sense; it is just storage. Both
uses are legitimate and they share an ancestor, but do not assume a coding
``tensor'' obeys the transformation rules until someone has told you what its
indices \emph{mean}. The data type is more permissive than the mathematical object.
\end{trapbox}

\section{Functional analysis: linear algebra in infinitely many dimensions}
\label{ch37:functional}

Every space in this book has been $\R^n$ or $\C^n$ for
some finite $n$. But the abstract definition of a vector space -- a set where you
can add and scale, obeying the usual rules -- never mentioned a dimension. What
happens if the vectors are \emph{functions}, and the space is infinite-dimensional?
Almost everything you know survives, with one subtle and important repair.

\subsection{Functions are vectors}
\label{ch37:functions-are-vectors}

Consider all continuous functions on the interval $[0,1]$. You can add two of them
pointwise, $(f+g)(x)=f(x)+g(x)$, and scale one, $(cf)(x)=cf(x)$, and every vector
space axiom holds. So this set of functions \emph{is} a vector space. But what is
its dimension? There is no finite list of functions whose combinations give every
continuous function. The space is infinite-dimensional.

\begin{intuitionbox}
Think of a function on $[0,1]$ as a vector with infinitely many components -- one
component $f(x)$ for each point $x$. A vector in $\R^3$ has three slots; a function
has a continuum of slots. The dot product of two vectors, $\sum_i u_i v_i$, becomes
a \emph{sum over all the slots}, which is an \emph{integral}:
\[
  \inner{f}{g} = \int_0^1 f(x)\,g(x)\,dx.
\]
Read the integral sign as ``add up the products component by component, where the
components are indexed by $x$ instead of by $i$.'' The inner product keeps the
same algebraic rules, with an integral in place of a finite sum.
\end{intuitionbox}

With an inner product comes \emph{everything geometric}. The length of a function
is $\norm{f} = \sqrt{\inner{f}{f}} = \sqrt{\int_0^1 f^2}$. Two functions are
\emph{orthogonal} when $\int_0^1 fg = 0$. The angle between functions, the
projection of one function onto another, the Cauchy--Schwarz inequality -- all of it
carries over verbatim, because every one of those facts was proved from the inner
product axioms alone, never from the dimension being finite.

\begin{example}[Two orthogonal functions]
\label{ch37:ex-orthfun}
On $[0,2\pi]$ with $\inner{f}{g}=\int_0^{2\pi} fg\,dx$, are $\sin x$ and $\cos x$
orthogonal? Compute:
\[
  \inner{\sin}{\cos} = \int_0^{2\pi}\sin x\cos x\,dx
  = \tfrac12\int_0^{2\pi}\sin 2x\,dx
  = \tfrac12\Big[-\tfrac12\cos 2x\Big]_0^{2\pi} = 0.
\]
Yes -- $\sin$ and $\cos$ are perpendicular vectors in function space. More is true:
the functions $\sin x$, $\cos x$, $\sin 2x$, $\cos 2x$, $\sin 3x$, and so on are
\emph{all} mutually orthogonal. You can check, for integers $m\neq n$, that $\int_0^{2\pi}\sin(mx)\sin(nx)\,dx=0$.
We are staring at an infinite orthogonal set -- an orthogonal \emph{basis} for a
space we cannot draw.
\end{example}

\begin{recoverybox}
\textbf{Pause and think.} What is the visual picture of $\int_0^{2\pi}\sin x\cos x\,dx = 0$? The product function $(\sin x)(\cos x) = \frac{1}{2}\sin 2x$ oscillates above and below the horizontal axis. Over one full cycle $[0,2\pi]$, the positive area on $[0,\pi/2]$ and $[\pi, 3\pi/2]$ is perfectly balanced by the negative area on $[\pi/2, \pi]$ and $[3\pi/2, 2\pi]$. Just as two vectors in $\R^2$ are orthogonal when their positive and negative coordinate products cancel ($u_1 v_1 + u_2 v_2 = 0$), two functions are orthogonal when their continuous product cancels across the domain.
\end{recoverybox}

\begin{figure}[H]
\centering
\begin{tikzpicture}[>=Latex]
\begin{axis}[width=12cm,height=5cm,
  axis lines=middle, xlabel={$x$},
  xtick={0,1.5708,3.1416,4.7124,6.2832},
  xticklabels={$0$,$\frac{\pi}{2}$,$\pi$,$\frac{3\pi}{2}$,$2\pi$},
  ytick={-1,1}, ymin=-1.4,ymax=1.4, domain=0:6.2832, samples=120,
  legend pos=outer north east, legend style={font=\small}]
  \addplot[vecblue,thick] {sin(deg(x))}; \addlegendentry{$\sin x$}
  \addplot[vecred,thick] {cos(deg(x))}; \addlegendentry{$\cos x$}
  \addplot[targetgreen,thick,dashed] {sin(deg(2*x))}; \addlegendentry{$\sin 2x$}
\end{axis}
\end{tikzpicture}
\caption{Three members of an infinite orthogonal family. Under the inner product
$\int_0^{2\pi}fg\,dx$, any two distinct curves here are perpendicular -- their
products integrate to zero. This is the basis behind Fourier series.}
\label{ch37:fig-sincos}
\end{figure}

\subsection{Fourier series: expanding in the orthonormal basis you cannot see}
\label{ch37:fourier}

In $\R^n$, once you have an
orthonormal basis $\vec{q}_1,\dots,\vec{q}_n$, you expand any vector by projecting
onto each basis vector and adding up:
\[
  \vec{v} = \sum_{k} \inner{\vec{q}_k}{\vec{v}}\,\vec{q}_k.
\]
Each coefficient $\inner{\vec{q}_k}{\vec{v}}$ is ``how much of $\vec{v}$ points along
$\vec{q}_k$.'' That expansion relies only on orthogonality, whether the sum has
three terms or infinitely many. Apply it to function space with
the trig basis, and out drops the \emph{Fourier series}:
\[
  f(x) = \frac{a_0}{2} + \sum_{k=1}^{\infty}\Big(a_k\cos kx + b_k\sin kx\Big),
  \qquad
  a_k = \frac{1}{\pi}\int_0^{2\pi} f(x)\cos kx\,dx,\quad
  b_k = \frac{1}{\pi}\int_0^{2\pi} f(x)\sin kx\,dx.
\]

\begin{focusbox}
Those integral formulas for $a_k$ and $b_k$ are standard orthogonal projections.
They follow directly from the projection formula
$\frac{\inner{\vec{q}_k}{\vec{v}}}{\norm{\vec{q}_k}^2}\,\vec{q}_k$
from the orthogonality chapter. Why is there a $1/\pi$ in front of the integrals,
and why is the constant term written as $a_0/2$? Because our trigonometric basis functions
are orthogonal, but not normalized to unit length:
\[
  \norm{\cos kx}^2 = \int_0^{2\pi}\cos^2(kx)\,dx = \pi, \qquad
  \norm{\sin kx}^2 = \int_0^{2\pi}\sin^2(kx)\,dx = \pi
\]
for every $k\ge 1$. Dividing by $\norm{\vec{q}_k}^2 = \pi$ produces the factor $1/\pi$.
For the constant basis function $q_0(x)=1$, its squared length is
$\norm{1}^2 = \int_0^{2\pi}1^2\,dx = 2\pi$, so its projection is
$\frac{1}{2\pi}\int_0^{2\pi}f(x)\,dx = \frac{a_0}{2}$. A Fourier coefficient is
literally the projection of a function onto an orthogonal basis direction. Fourier
analysis \emph{is} the ``expand in an orthogonal basis'' theorem, run in infinite dimensions.
\end{focusbox}

Let us see one coefficient computed, so the abstraction touches the ground.

\begin{example}[A Fourier coefficient as a projection]
\label{ch37:ex-fourier}
Take the simple function $f(x)=x$ on $[0,2\pi]$ and project it onto $\sin x$. The
coefficient is
\[
  b_1 = \frac{1}{\pi}\int_0^{2\pi} x\sin x\,dx.
\]
Integrate by parts with $u=x$, $dv=\sin x\,dx$, so $du=dx$, $v=-\cos x$:
\[
  \int_0^{2\pi} x\sin x\,dx
  = \big[-x\cos x\big]_0^{2\pi} + \int_0^{2\pi}\cos x\,dx
  = -2\pi\cos 2\pi + \big[\sin x\big]_0^{2\pi}
  = -2\pi + 0 = -2\pi.
\]
Therefore $b_1 = \frac{1}{\pi}(-2\pi) = -2$. The function $f(x)=x$ has component $-2$
along $\sin x$, directly analogous to how $\colvec{3\\4}$ has component $3$ along $\vec{e}_1$.
Same operation, infinite-dimensional setting.
\end{example}

\begin{strangbox}
Strang devotes a beautiful section to the Fourier matrix and the Discrete Fourier
Transform -- the \emph{finite}, $n$-point version, where the orthonormal basis is the
columns of a unitary matrix and the ``integral'' is a plain finite sum. The continuous
Fourier series of this section is what the DFT becomes as the number of sample
points goes to infinity. Read both together and the message is unmistakable: the DFT
is a change of basis to an orthonormal basis (Strang's emphasis), and the Fourier
series is the same change of basis when the vectors are functions.
\end{strangbox}

\subsection{The repair: completeness and Hilbert space}
\label{ch37:hilbert}

Infinite dimensions introduce one important subtlety that never arises in $\R^n$. In $\R^n$,
every \emph{Cauchy sequence} (a sequence whose terms bunch closer and closer together,
so $\norm{\vec{v}_m-\vec{v}_n}\to 0$ as $m,n\to\infty$) automatically converges to a
limit that lives safely inside $\R^n$. In finite dimensions, you can never escape the
space simply by taking limits.

In infinite dimensions, this safety net snaps. You can easily build a sequence of
smooth, continuous functions that get closer and closer to one another under the
$L^2$ inner-product distance, yet whose limit develops a vertical cliff -- a jump
discontinuity. Figure~\ref{ch37:fig-completeness} shows how this happens: the ramp
steepens, the difference area between successive approximations vanishes, yet the limit
function is a discontinuous step function $\operatorname{sgn}(x)$, which has been expelled
from the space of continuous functions $C[-1,1]$. A space where this defect never
happens -- where every Cauchy sequence has its limit safely inside -- is called
\emph{complete}. An infinite-dimensional inner-product space that is complete is a
\emph{Hilbert space}.

\begin{figure}[H]
\centering
\begin{tikzpicture}[scale=1.2,>=Latex]
  \draw[->,thick,spacegray] (-2.5,0)--(2.7,0) node[right,font=\scriptsize]{$x$};
  \draw[->,thick,spacegray] (0,-1.5)--(0,1.6) node[above,font=\scriptsize]{$y$};

  % Ticks
  \draw (-2,0.06)--(-2,-0.06) node[below=3pt,font=\scriptsize] {$-1$};
  \draw (2,0.06)--(2,-0.06) node[below=3pt,font=\scriptsize] {$1$};
  \draw (0.06,1)--(-0.06,1) node[left=3pt,font=\scriptsize] {$1$};
  \draw (0.06,-1)--(-0.06,-1) node[right=3pt,font=\scriptsize] {$-1$};

  % Discontinuous limit step function
  \draw[gray!70,very thick,dashed] (-2,-1)--(0,-1);
  \draw[gray!70,very thick,dashed] (0,1)--(2,1);
  \draw[gray!50,dotted,thick] (0,-1)--(0,1);
  \node[font=\scriptsize,text=spacegray!90!black,anchor=west] at (0.3,1.3) {step limit $\operatorname{sgn}(x)$};

  % Approximations f_n
  \draw[sourceblue,thick] (-2,-1)--(-1,-1)--(1,1)--(2,1);
  \draw[targetgreen!80!black,thick] (-2,-1)--(-0.5,-1)--(0.5,1)--(2,1);
  \draw[vecred,thick] (-2,-1)--(-0.25,-1)--(0.25,1)--(2,1);

  % Legend in upper-left corner
  \node[draw=spacegray!40,fill=white,rounded corners=2pt,font=\tiny,anchor=north west,inner sep=3pt,align=left] at (-2.4,1.45) {
    \textcolor{sourceblue}{\textbf{---}} $f_1$\\
    \textcolor{targetgreen!80!black}{\textbf{---}} $f_2$\\
    \textcolor{vecred}{\textbf{---}} $f_4$
  };

  \fill[spacegray] (0,0) circle (1.5pt);
\end{tikzpicture}
\caption{Why the space of continuous functions $C[-1,1]$ fails to be complete. The continuous functions $f_n(x)$ ramp linearly from $-1$ to $+1$ over the narrowing window $[-1/n, 1/n]$. As $n\to\infty$, the difference between $f_m$ and $f_n$ has shrinking area, so $(f_n)$ is a Cauchy sequence in the $L^2$ norm. Yet its limit is the discontinuous step function $\operatorname{sgn}(x)$, which is not continuous. A Hilbert space repairs this defect by including all such limits.}
\label{ch37:fig-completeness}
\end{figure}

\begin{formalbox}
A \emph{Hilbert space} is a vector space $H$ with an inner product $\inner{\cdot}{\cdot}$
that is \emph{complete}: every sequence $(\vec{v}_n)$ with $\norm{\vec{v}_m-\vec{v}_n}\to 0$
as $m,n\to\infty$ converges to some limit \emph{in} $H$. Finite-dimensional inner-product
spaces ($\R^n$, $\C^n$) are automatically Hilbert spaces. The fundamental
infinite-dimensional example is $L^2[a,b]$, the square-integrable functions with
$\inner{f}{g}=\int_a^b \overline{f}g\,dx$.
\end{formalbox}

Why does completeness matter? Because it is exactly what makes the infinite Fourier
sum \emph{mean something}. Each partial sum of a Fourier series is a function in our
space; completeness guarantees the limit is too, so ``$f$ equals its Fourier series''
is a legitimate equation (in the $L^2$ sense of length). Completeness is the essential
property required for infinite series to have well-defined limits in the space, which is
why Hilbert spaces form the standard setting for functional analysis.
Inside a Hilbert space, the entire geometric apparatus of this book -- bases,
projections onto subspaces, the Pythagorean theorem, even a version of the spectral
theorem for the right operators -- holds exactly as you would hope.

\begin{aestheticbox}
The geometric framework built from vectors in $\R^n$ -- length, angle,
perpendicularity, projection, and orthonormal expansion -- extends naturally to
infinite-dimensional spaces. Quantum mechanics is formulated in a Hilbert space,
where observables are self-adjoint operators, states are unit vectors, and
measurable quantities are eigenvalues. Similarly, the heat and wave equations are
solved by expanding in an orthonormal basis of eigenfunctions, decoupling the
system across modes. The projection formula applies equally to finite-dimensional
vectors and continuous functions, serving as a unifying foundation for signal
processing, quantum physics, and differential equations.
\end{aestheticbox}

\begin{connectionbox}
\textbf{Operators as infinite matrices.} In a Hilbert space, the role of a matrix is
played by a \emph{linear operator} on functions -- differentiation $\frac{d}{dx}$,
integration, multiplication by $x$. Many of these have eigenfunctions and
eigenvalues. The operator $-\frac{d^2}{dx^2}$ on $[0,\pi]$ (with the function
vanishing at the ends) has eigenfunctions $\sin(kx)$ with eigenvalues $k^2$ -- a
``diagonalization'' with infinitely many eigenpairs. Solving a differential equation
by Fourier series \emph{is} eigenvector decomposition for this operator. The
spectral theorem for symmetric matrices, $\mat{A}=\sum\lambda_k\vec{q}_k\vec{q}_k\T$,
has a Hilbert-space sibling that powers the entire theory of partial differential
equations.
\end{connectionbox}

\section{Looking back, looking out}
\label{ch37:lookback}

Stand back and look at the three doors together, because they are really one
doorway seen from three angles.

\begin{itemize}
  \item \textbf{Groups} took the matrices that preserve geometry and asked about
        \emph{their} structure -- how they combine, what symmetries they encode,
        and how a symmetry can \emph{act} on a space. A central theme was
        simultaneous diagonalization: the irreducible representations correspond
        to the simultaneous eigen-decomposition of a symmetry group.
  \item \textbf{Tensors} took linearity-in-one-slot and let it happen in several
        slots at once. Outer products, the determinant as an alternating form, and
        the multi-index arrays of machine learning are all the same idea: stack the
        bilinear structure you already had, index by index.
  \item \textbf{Functional analysis} took the dimension to infinity and discovered
        that geometry -- orthogonality, projection, bases -- not only survives but
        becomes the natural language of Fourier analysis, differential equations,
        and quantum mechanics, once you pay the price of completeness.
\end{itemize}

\begin{walkbox}
Take a walk and try to find, in your own life or studies, one example behind each
door. A symmetry of an object you can pick up (a group). A quantity that depends
linearly on two separate things (a tensor). A signal -- music, a heartbeat, a
spoken word -- that is naturally a sum of pure tones (a Fourier expansion in a
function space). When you can name one of each from memory, the abstract chapter
above has become yours.
\end{walkbox}

None of this replaces the careful study these subjects deserve; each door opens onto
a long hallway.

\begin{focusbox}
The conceptual arc of the book: vectors and span define a \emph{space};
the dot product defines \emph{geometry}; matrices represent \emph{linear maps};
eigenvalues and the SVD identify \emph{optimal coordinates}. This chapter showed
these principles extending to groups, tensors, and infinite dimensions.
\end{focusbox}

\subsection*{Boundaries in multilinear and infinite-dimensional spaces}

Venturing into these wider territories does come with honest boundaries. In tensor
spaces, dimensions multiply exponentially ($\dim V^{\otimes k} = d^k$), so tracking
just 50 interacting quantum particles requires $2^{50}$ amplitudes---exceeding the
memory of supercomputers and forcing physicists to invent tensor network approximations.
In infinite-dimensional Hilbert spaces, closed unit spheres are no longer compact,
operators can have continuous spectra without a single discrete eigenvector (like
position and momentum in quantum physics), and differential operators like
$\frac{d}{dx}$ are unbounded, requiring careful topological care with dense domains
and Sobolev spaces.

For reference, two companion appendices provide background: \textbf{Appendix~A}
reviews sets, mappings, and field axioms, and \textbf{Appendix~B} reviews summation
notation, complex arithmetic, and matrix calculus.

\section*{Exercises}
\addcontentsline{toc}{section}{Exercises}

\begin{exercise}[Mental check: Determinants of orthogonal matrices]
\label{ch37:ex-mental-det}
Without writing down any coordinates, use the identity $\mat{Q}\T\mat{Q} = \mat{I}$ and
the determinant product rule $\det(\mat{A}\mat{B}) = (\det\mat{A})(\det\mat{B})$ to
show in one line that every orthogonal matrix must have $\det\mat{Q} = \pm 1$.
Can an orthogonal matrix have $\det\mat{Q} = 2$? Why geometrically can a rigid motion
never scale volumes?
\end{exercise}

\begin{exercise}[Closure of $SO(2)$ by computation]
\label{ch37:ex-so2}
Multiply $\mat{R}(30^\circ)$ and $\mat{R}(45^\circ)$ as $2\times 2$ matrices and
verify directly that the product equals $\mat{R}(75^\circ)$. Use your computation to
state the angle-addition formula for $\cos(\alpha+\beta)$ that appears in the
top-left entry.
\end{exercise}

\begin{exercise}[Reflections are orthogonal but not rotations]
\label{ch37:ex-reflection}
Let $\mat{F}=\begin{bmatrix}0&1\\1&0\end{bmatrix}$. Verify $\mat{F}\T\mat{F}=\mat{I}$
(so $\mat{F}\in O(2)$) and compute $\det\mat{F}$. To which half of $O(2)$ does
$\mat{F}$ belong? Find its eigenvalues and eigenvectors and describe the reflection
geometrically (across which line?).
\end{exercise}

\begin{exercise}[Euler's axis in action]
\label{ch37:ex-euler}
Consider the rotation of $\R^3$ that cycles the axes, $\mat{R}\colvec{x\\y\\z}
=\colvec{z\\x\\y}$, with matrix
$\mat{R}=\begin{bmatrix}0&0&1\\1&0&0\\0&1&0\end{bmatrix}$. Verify $\mat{R}\in SO(3)$
(orthogonal, determinant $+1$). Then find a nonzero vector with $\mat{R}\vec{x}
=\vec{x}$ -- the axis guaranteed by Euler's theorem.
\end{exercise}

\begin{exercise}[A two-element representation]
\label{ch37:ex-reptwo}
The cyclic group of order $2$ can be represented on $\R^2$ by sending the nontrivial
element to $\mat{A}=\begin{bmatrix}-1&0\\0&1\end{bmatrix}$. Check that $\mat{A}^2
=\mat{I}$ (so the group law is respected). Is this representation already split into
one-dimensional invariant pieces? Identify the two invariant lines and say what the
group element does on each.
\end{exercise}

\begin{exercise}[Outer products and rank]
\label{ch37:ex-outer}
Let $\vec{u}=\colvec{2\\-1\\3}$ and $\vec{v}=\colvec{1\\4}$. Compute the outer
product $\vec{u}\,\vec{v}\T$ as a $3\times 2$ matrix. Without further computation,
state its rank and explain why every nonzero $3\times 2$ outer product has that same rank.
\end{exercise}

\begin{exercise}[The determinant is alternating]
\label{ch37:ex-altdet}
Using only multilinearity and the alternating property (sign flips when two columns
are swapped), prove that a $2\times 2$ matrix with two equal columns has determinant
$0$. Then explain in one sentence why this matches the geometric meaning of
determinant as signed area.
\end{exercise}

\begin{exercise}[Orthogonality of functions]
\label{ch37:ex-orthfun2}
On $[0,2\pi]$ with $\inner{f}{g}=\int_0^{2\pi}fg\,dx$, verify that $\sin x$ and
$\sin 2x$ are orthogonal. (Hint: use the product-to-sum identity $\sin A\sin B
=\tfrac12[\cos(A-B)-\cos(A+B)]$.) What is $\norm{\sin x}^2=\int_0^{2\pi}\sin^2 x\,dx$?
\end{exercise}

\begin{exercise}[A Fourier coefficient]
\label{ch37:ex-fouriercoef}
For $f(x)=x$ on $[0,2\pi]$, compute the coefficient $a_1=\frac1\pi\int_0^{2\pi}
x\cos x\,dx$ by integration by parts. Combined with the worked example $b_1=-2$,
write the first sine-and-cosine terms of the Fourier series of $f$.
\end{exercise}

\begin{exercise}[Projection in function space]
\label{ch37:ex-funproj}
In $\R^n$ the projection of $\vec{v}$ onto a unit vector $\vec{q}$ is
$\inner{\vec{q}}{\vec{v}}\vec{q}$. Translate this into function space: write the
formula for the projection of a function $f$ on $[0,2\pi]$ onto the (not-yet-unit)
direction $\sin x$, remembering to divide by $\norm{\sin x}^2$. This single term is
the ``best multiple of $\sin x$'' approximating $f$ in the least-squares sense.
\end{exercise}

\begin{exercise}[Why completeness matters -- a discussion problem]
\label{ch37:ex-complete}
Consider continuous functions on $[-1,1]$ and a sequence $f_n$ that is $-1$ for
$x<-\tfrac1n$, $+1$ for $x>\tfrac1n$, and linear in between. Sketch $f_1,f_2,f_3$.
Argue informally that the $f_n$ get closer and closer to one another in the
$L^2$ length $\norm{f_m-f_n}$, yet their ``limit'' is a discontinuous step function.
Explain in a sentence why this shows the continuous functions are \emph{not}
complete, motivating the larger Hilbert space $L^2$.
\end{exercise}

\begin{exercise}[Challenge: the trace is a tensor contraction]
\label{ch37:ex-trace}
The trace of a matrix is $\tr\mat{A}=\sum_i A_{ii}$. Viewing $\mat{A}$ as an
order-2 tensor $A_{ij}$, the trace ``contracts'' the two indices by setting them
equal and summing. Prove the cyclic property $\tr(\mat{A}\mat{B})=\tr(\mat{B}\mat{A})$
by writing both sides as the double sum $\sum_{i,j}A_{ij}B_{ji}$, and explain why
this makes the trace independent of the basis (i.e. $\tr(\mat{P}\inv\mat{A}\mat{P})
=\tr\mat{A}$).
\end{exercise}

\begin{exercise}[Challenge: $U(1)$ and the circle]
\label{ch37:ex-u1}
Show that the map $\theta\mapsto e^{i\theta}$ from angles to unit complex numbers
sends addition of angles to multiplication of complex numbers, and that
$\theta\mapsto\mat{R}(\theta)$ does the same into $SO(2)$. Conclude that the unit
circle, the group $SO(2)$, and the unit complex numbers $U(1)$ are ``the same group''
(isomorphic). What is the inverse of $e^{i\theta}$, and what rotation does it
correspond to?
\end{exercise}
